\chapter*{Conclusion}
\markboth{Conclusion}{}
\addstarredchapter{Conclusion}

Ce projet a permis de construire, de bout en bout, une application capable de lire automatiquement les informations d'un document d'identité (carte nationale d'identité, passeport, certificat de scolarité) et de les comparer à une base de référence pour détecter d'éventuelles incohérences. Le système combine du traitement d'image (OpenCV), de la reconnaissance optique de caractères (Tesseract), et de la comparaison floue de chaînes de caractères (rapidfuzz), assemblés derrière une interface web simple (Streamlit) et livrés dans un conteneur Docker.

L'ensemble des besoins exprimés au début du projet a été couvert : extraction sur quatre types de documents (carte nationale d'identité, passeport, certificat de scolarité, titre de séjour), comparaison avec détection d'incohérences, interface utilisateur, possibilité de charger sa propre base de référence, fonctionnement en français et en conteneur. Les 41 tests automatisés passent à 100\,\%.

Le point le plus marquant de ce projet n'est cependant pas la première version fonctionnelle, mais le travail mené ensuite pour en mesurer et en améliorer la précision, en deux temps. Un premier travail, mené sur les trois types de documents initiaux et guidé par un banc d'essai chiffré (chapitre~5), a révélé trois défauts cachés (un risque de blocage de l'OCR sans limite de temps, une erreur de signe qui doublait l'inclinaison des images au lieu de la corriger, et une mauvaise hypothèse sur la convention d'angle d'une fonction d'OpenCV), qui faisaient chuter fortement la précision sur des cas pourtant courants ; leur correction a fait passer la précision globale mesurée de 57,1\,\% à 79,1\,\%, soit un gain de 22 points. Un second travail, mené lors de l'ajout du titre de séjour comme quatrième type de document et validé cette fois sur de vraies photographies plutôt que sur le banc d'essai synthétique, a révélé cinq défauts supplémentaires détaillés au chapitre~4 : un ordre de lecture du texte infidèle à la mise en page réelle, un fond de sécurité imprimé rendant le texte illisible, une collision entre deux champs récupérant par erreur la même valeur, un format de carte d'identité récent non reconnu, et une zone MRZ de passeport pouvant mélanger deux lectures indépendantes. La leçon que je retiens des deux temps de ce travail est qu'une architecture propre et un code qui semble correct à la lecture ne garantissent pas un bon comportement : seule une mesure objective — chiffrée quand c'est possible, ou à défaut un test manuel systématique sur des cas réels — permet de savoir si un système fonctionne réellement bien, et de prouver qu'une correction a un effet réel plutôt que supposé.

Le système réalisé garde malgré tout des limites claires. La précision reste imparfaite sur les images de très basse résolution, où un détail déjà perdu ne peut pas toujours être récupéré par un simple agrandissement ; une correction de perspective (pour les photos prises avec un angle de vue, et non plus seulement une rotation dans le plan de l'image) pourrait améliorer ce point. Le banc d'essai ne teste qu'un seul sens de rotation alors que la correction des deux sens a été validée séparément par test unitaire ; il serait utile de l'étendre. Le titre de séjour, quatrième type de document ajouté après la mise en place du banc d'essai, n'y est pour l'instant pas intégré : sa validation repose uniquement sur des tests manuels, ce qui est moins rigoureux et moins reproductible qu'une mesure chiffrée ; l'étendre à ce nouveau type de document, ainsi qu'écrire des tests automatisés couvrant les cinq défauts corrigés en fin de chapitre~4, est la prochaine étape naturelle. Le fichier Excel utilisé comme base de référence resterait par ailleurs à remplacer par une vraie base de données pour un usage à plus grande échelle. Enfin, comme il s'agit de données d'identité, une utilisation réelle demanderait des garanties de sécurité (chiffrement, suppression des images après traitement) que ce projet n'a pas traitées : l'objectif ici était la chaîne de traitement OCR elle-même.

Ces limites ouvrent autant de pistes pour la suite : ce projet constitue une base solide et testée, sur laquelle il serait possible de continuer à travailler, que ce soit en ajoutant de nouveaux types de documents, en renforçant la sécurité, ou en poursuivant le travail de mesure et d'amélioration de la précision engagé ici.
